% Source: https://artofproblemsolving.com/community/c3h2788839_2009_manhattan_mathematical_olympiad_grades_78_problems_14
% Stable identity: aops-topic-2788839-q02
Use the second piece both as a straightedge and as a template. Call the first piece $ABCD$ with $AB=73$ and $BC=19$.

Lay the second piece on the first so that two opposite corners of the second piece sit on $A$ and $C$ (possible because the two pieces have equal diagonals; the second piece is then the mirror image of the first in the line $AC$). Trace its long edges: they cross $AB$ at $P$ and $CD$ at $Q$. The overlap $APCQ$ is a rhombus (two strips of equal width), so its diagonals $AC$ and $PQ$ bisect each other, and $PQ$ passes through the midpoint of $AC$, the centre of $ABCD$. Draw $PQ$ using an edge of the second piece ($PQ\approx19.6$ cm, short enough).

Repeat with the opposite corners $B$ and $D$ to get the segment $P'Q'$ (the mirror image of $PQ$). The centre is where $PQ$ and $P'Q'$ cross.
\begin{asy}
usepackage("amsmath");
size(7cm);
defaultpen(fontsize(10pt));
pair A=(0,0),B=(73,0),C=(73,19),D=(0,19);real s=38.97;pair P=(s,0),Q=(73-s,19),P2=(73-s,0),Q2=(s,19);
real t=2*atan2(19,73);pair u=(Cos(t*180/pi),Sin(t*180/pi)),n=(u.y,-u.x);
draw(A--B--C--D--cycle);draw(A--A+75.43*(Cos(atan2(19,73)*180/pi),Sin(atan2(19,73)*180/pi)),dotted);
draw(A--A+(73*u)--A+(73*u)+19*n--A+19*n--cycle,dashed);draw(P--Q,linewidth(1));draw(P2--Q2,linewidth(1));dot((36.5,9.5));
label("$A$",A,SW);label("$B$",B,SE);label("$C$",C,NE);label("$D$",D,NW);label("$P$",P,S);label("$Q$",Q,N);label("$P'$",P2,S);label("$Q'$",Q2,N);
\end{asy}

(The diagonal $AC\approx75.4$ cm is longer than any edge, which is why it is not drawn directly.)
